\chapter{如何应对紧急的交期} \label{ch:deadline-ai}

前面各章讨论的都是“如何把代码写好”：状态机、任务调度、串口协议、菜单系统……这些知识解决的是“会不会做”的问题。但当项目真正进入倒计时，交期只剩一周甚至一个晚上时，真正决定成败的往往是另一类能力：如何调度自己的时间、如何调度别人的时间、如何调度 AI 这个“新员工”。本章不介绍具体代码，而是结合作者在第~\ref{ch:about-author}~章里讲述的真实经历，讨论校园和职场两种场景下的紧急交期处理。本章的核心观点只有三条：
\begin{enumerate}
\item AI 是廉价的“体力劳工”，不是免费的“脑力军师”，决策权永远不能交给它；
\item 越是时间不够，越要用确定性工具，避开看不懂的“魔法”和养不活的“轮子”；
\item 无论是同学、老师还是老板，谈判的本质都是把“必须交付什么”说清楚，而不是比拼谁的代码写得快。
\end{enumerate}

% ---------------------------------------------------------------------
\section{使用 AI Agent 高效编程}

AI Agent 和普通聊天 AI 最大的区别是：它不止“回答你”，还会真的打开你的工程、读代码、改文件、运行命令。对嵌入式开发来说，这就相当于多了一个不知疲倦、手速极快、但偶尔会自作聪明的实习生。用得好，它是交期压力下最强的外援；用不好，它会把一个能跑的项目改成一个“看起来能跑”的项目。

% ---------------------------------------------------------------------
\subsection{DeepSeek Codex 的部署}

Codex 是 OpenAI 推出的 AI Agent 编程平台，它的工作方式不是“你在网页上问一句，它贴一段代码”，而是把整个项目目录当作工作区：先读懂工程结构和现有代码，再按任务修改文件、运行构建命令，最后把改动结果交回给你确认。

作者的习惯是：把 AI Agent 当成“自动代码补全器 + 体力劳工”。生成格式化代码、填充注释、写测试桩、批量重构这类重复劳动可以放心交给它；但“用什么方案、砍哪个功能、代码要不要保留”这类决定，永远自己拍板。就像第~\ref{ch:about-author}~章说的：Agent 可以替你写尽天下代码，但不能替你按下那枚“锤子”。

% ---------------------------------------------------------------------
\subsection{向 DeepSeek 国产大模型开发平台付费}

没有免费的午餐。Codex 虽然是 OpenAI 推出的 AI Agent 编程平台，但国内无法访问 ChatGPT 的任何平台，因此我们需要向国产大模型开发平台付费购买“AI 体力劳动”。这里的“付费”可以理解成两笔账：

\begin{enumerate}
\item 给平台的 Token 费：DeepSeek 等国产平台提供与 OpenAI 兼容的 API，注册后充值、拿到密钥，再把 Codex 的模型地址指向该平台即可使用。费用通常按 Token 计费，生成一段代码的开销很小，但比“0 元试用”稳定得多；
\item 给 AI 的“工时费”：作者在第~\ref{ch:about-author}~章提到“公司 Token = 免费算力，往死里用”，本质上是把低认知密度的体力劳动外包给 Token，把自己省下来的时间花在真正需要思考的地方。
\end{enumerate}

值得强调的是：付费买来的是“体力”，不是“决策权”。AI 生成 100 行代码只需要几秒钟，但如果它把架构方向带偏了，后面可能要花 100 个小时来还债。所以每次让 AI 动手前，先把边界说清楚：改哪个文件、遵守什么命名规范、不许动哪些核心逻辑。

% ---------------------------------------------------------------------
\subsection{解决物理世界与虚拟世界的调试麻烦}

嵌入式开发和纯软件开发的本质区别，是它同时存在于两个世界：

\begin{itemize}
\item 虚拟世界：源码、编译器、仿真器、串口日志，Agent可以自由进出；
\item 物理世界：单片机引脚、电平、示波器波形、电机转动、屏幕刷新，Agent永远只能“听说”不能“看见”。
\end{itemize}

很多初学者让 AI 改完代码后直接下载到板子上，发现不行，又让 AI 继续瞎猜，结果越改越乱。正确做法是：让 AI 在虚拟世界里工作，把物理世界的结果“翻译”给它。

\begin{enumerate}
\item 先在 PC 端搭仿真（例如本书菜单章节里的 EasyX 仿真、串口助手、命令行测试），让 AI 能在虚拟世界里反复编译运行；
\item 每次真机调试只改一个变量，把示波器截图、串口日志、打印信息原样贴给 AI；
\item 编译、下载、点“运行”这个最后的动作，尽量由人亲手完成。
\end{enumerate}

作者把编译报错称为“理解系统的地震波”，把链接成功称为“亲手点燃的火药”。原因就在这里：物理世界的现象是判断 AI 是否说谎的唯一标准。AI 可以在虚拟世界里撒谎，但示波器不会。

% ---------------------------------------------------------------------
\subsection{\texttt{AGENTS.md} 的参考声明}

\texttt{AGENTS.md} 是放在项目根目录（或子目录）里的说明文件，专门用来告诉 AI Agent：这个项目是怎么组织的、代码要按什么规范写、哪些事情绝对不能做。它相当于给 AI 发的 “员工手册”。

以本仓库的 \texttt{AGENTS.md} 为例，里面就写明了：

\begin{itemize}
\item 所有代码注释、配置说明、文档备注一律使用中文；
\item 函数必须有文档注释，结构体、枚举、命名必须遵守项目命名规范；
\item 代码注释密度不宜过密，但每个代码组都要有说明；
\item 指针判空必须写 \texttt{if (ptr != NULL)}，不能直接写 \texttt{if (ptr)}；
\item 生成临时文件后必须主动清理。
\end{itemize}

在紧急交期下，\texttt{AGENTS.md} 的价值不是“显得专业”，而是省时间：

\begin{enumerate}
\item AI 不会反复写出风格不一致的代码，你不用每次重写大段注释和命名；
\item 换一个 AI 或换一个会话，行为规范仍然一致；
\item 作者还可以在 \texttt{AGENTS.md} 里写下“最终防线”，例如“不要替我编译，把编译结果交给我看”，把 AI 限制在安全的工作范围内。
\end{enumerate}

注意区分两类内容：只对当前这一次任务有效的指令，写在对话提示词里就够了；希望所有会话、所有 Agent 都遵守的长期规则，才值得写进 \texttt{AGENTS.md}。下面给出一份简化版的 \texttt{AGENTS.md} 参考示例：

\begin{lstlisting}
# AGENTS.md
> 本文件用于配置 AI 助手的行为规范。

## 全局指令
- 语气：友好、耐心、清晰，尽量使用中文
- 注释语言：所有代码注释、配置说明、文档备注一律使用中文

## 命名规范
- 函数：模块名_功能名（动宾短语），模块层级不超过 2 层
- 结构体 / 枚举：大驼峰命名 + _T 后缀
- 宏定义：全大写下划线
- 函数指针：以 Callback 结尾

## 代码要求
- 每个函数必须有文档注释，说明作用、参数和返回值
- 每个代码组都要有注释，组与组之间空行分隔
- 指针判空必须写 if (ptr != NULL)，不能直接写 if (ptr)

## 禁止事项
- 不要替我编译，把编译结果给我看
- 不要修改 core 目录下的协议定义和对外接口
- 不要删除我没有要求删除的文件
\end{lstlisting}

% ---------------------------------------------------------------------
\subsection{有必要写 \texttt{SKILL.md} 吗}

\texttt{SKILL.md} 用来描述“一项可以反复执行的工作流”。如果说 \texttt{AGENTS.md} 是 “员工手册”，那么 \texttt{SKILL.md} 就是“标准作业程序”。

举例来说，如果每次给菜单系统新增一个页面，都要依次完成以下步骤：创建菜单项结构体、注册回调函数、更新菜单栈、在 PC 仿真里验证、再移植到单片机，那么就可以把它写成一个 Skill。之后 AI 接到“新增一个页面”的需求时，就会自动按这个流程执行，而不是每次从零猜测。

判断标准很简单：

\begin{itemize}
\item 这个操作只做一两次，或每次做法都不一样 $\rightarrow$ 不写，直接在对话里说；
\item 这个操作会反复出现，且步骤基本固定 $\rightarrow$ 值得写；
\item 操作涉及硬件下载、编译命令、回归测试等不能出错的环节 $\rightarrow$ 强烈建议写。
\end{itemize}

作者的观点是：AI 既然是“体力劳工”，就更应该给它一份固定的“操作手册”，而不是每次都现场教。写 \texttt{SKILL.md} 表面上多花了半小时，但它能把“每次重复劳动都要人盯着”变成 “一次标准化，之后全部自动”，这正是第~\ref{ch:about-author}~章里“用确定性的劳动对冲不确定性风险”的翻版。下面给出一份简单的 \texttt{SKILL.md} 参考示例：

\begin{lstlisting}
# SKILL.md

## 技能名称
新增菜单页面

## 适用场景
需要给现有菜单系统新增一个功能页面时。

## 操作步骤
1. 在 menu_user.c 中创建菜单项结构体，并注册回调函数
2. 在菜单表中登记新页面，设置跳转入口
3. 在 PC 仿真中验证页面显示与按键逻辑
4. 移植到单片机，确认 LCD 与按键驱动正常
5. 运行回归测试，检查原有菜单页面没有受影响

## 注意事项
- 不要修改菜单核心文件 menu.c / menu.h
- 新增 API 必须补充中文文档注释
- 命名必须遵守 AGENTS.md 中的规范
\end{lstlisting}

% ---------------------------------------------------------------------
\subsection{有必要写 \texttt{MEMORY.md} 吗}

\texttt{MEMORY.md}（有的工具里叫 \texttt{MEMORYS.md}）用于保存跨越多个会话的长期信息，比如作者的个人偏好、常用工具链、曾经踩过的坑。它和 \texttt{AGENTS.md} 的分工是：

\begin{itemize}
\item \texttt{AGENTS.md}：面向“当前仓库”，告诉 AI 这个项目该怎么干活；
\item \texttt{MEMORY.md}：面向“长期记忆”，告诉 AI 这个用户习惯怎么干活。
\end{itemize}

是否有必要写，取决于你使用 AI 的频率：

\begin{enumerate}
\item 只是偶尔用一次 AI 写小程序 $\rightarrow$ 没必要，记忆文件反而会成为过时信息的来源；
\item 每天都在同一个开发环境里使用 AI Agent $\rightarrow$ 建议写，至少记录“中文注释、优先使用确定性工具、编译由人执行”这类稳定的个人偏好；
\item 跨项目使用同一套习惯 $\rightarrow$ 可以把通用偏好放记忆，把项目特有规范放 \texttt{AGENTS.md}。
\end{enumerate}

写 \texttt{MEMORY.md} 时要注意两点：第一，只记录“长期不变”的偏好，不要记录会过期的事实；第二，记忆必须定期清理。作者的原则是：AI 可以记住我的习惯，但不能用一条过时的记忆替我做决定。下面给出一份简单的 \texttt{MEMORY.md} 参考示例：

\begin{lstlisting}
# MEMORY.md

## 用户偏好
- 代码注释、文档说明一律使用中文
- 命名遵守仓库 AGENTS.md 的规范
- 优先使用确定性工具，慎用看不懂的黑盒框架
- AI 负责体力劳动，架构决策由用户完成
- 编译、下载、烧录等动作由用户亲手执行

## 常用技术栈
- 单片机：STM32 系列、STC 系列
- 语言：C / C++ / Python
- 工具：Keil、VS Code、Codex、Git

## 注意事项
- 记忆会过期，遇到冲突时以当前对话和 AGENTS.md 为准
\end{lstlisting}

% ---------------------------------------------------------------------
\subsection{其他 AI 常见术语及其含义}

下面整理交期场景下最常遇到的 AI 术语，方便读者阅读资料和向别人描述需求：

\begin{itemize}
\item Token：大模型计费和输入输出的最小单位，可以粗略理解成“字数”。写代码时把需求说清楚、把无关文件收起来，就是在省 Token；
\item Prompt 与 System Prompt：你给 AI 的指令。系统提示词决定 AI 的“人设”，普通对话决定本次任务内容；
\item 上下文窗口（Context Window）：AI 一次能“记住”的最大内容量。工程太大时不能把整个仓库塞进去，需要靠 \texttt{AGENTS.md}、索引或按目录分段喂；
\item Agent：能自主调用工具、读写文件、执行命令的 AI，而不是只回一句话的聊天机器人；
\item Function Call 与 Tool Use：AI 调用外部函数或工具的能力，例如读文件、执行编译、查串口日志，这是 Agent 能“干活”的基础；
\item MCP（Model Context Protocol）：统一 AI 与外部工具或数据源连接的协议，可以理解成 AI 的“USB 接口”，插上什么工具就能用什么能力；
\item RAG（检索增强生成）：先从资料库检索相关内容，再让 AI 基于检索结果回答，用来缓解“凭记忆瞎编”的问题；
\item 微调（Fine-tuning）：用专门数据进一步训练模型，成本高，普通项目几乎用不到，别把它和“写提示词”混为一谈；
\item 幻觉（Hallucination）：AI 一本正经地编造不存在的 API、寄存器或文件。判断方法只有一个：让它运行给你看；
\item 沙箱（Sandbox）：AI 只能在一个受限环境里执行命令，防止误删文件或破坏系统，是保护项目安全的重要机制。
\end{itemize}

面对这些术语，作者的建议始终是：术语只是工具的名字，重点是“谁能替我干活、谁能替我决定”。能看透底牌的工具就多用，看不透的“魔法”就少碰。

% ---------------------------------------------------------------------
\section{校园的紧急交期处理}

校园项目的交期通常来自课程设计、竞赛和毕业设计。它的特点是：时间短、组员水平参差不齐、硬件资源有限，而且很多时候“通过”的标准并不高——但没人告诉你标准到底是什么。

% ---------------------------------------------------------------------
\subsection{同学间的小组分工分配方法}

小组分工最常见的错误，是把任务按“代码量”平均分配：你写一半，我写一半。结果接口对不上、风格不统一、最后两天全组都在联调。

正确的思路是按“接口”和“风险”分工：

\begin{enumerate}
\item 先定接口，再分任务：例如菜单项目先约定好“菜单项结构体长什么样、按键事件怎么传递、界面刷新由谁调用”，之后各组员在自己的模块里互不干扰；
\item 让最靠谱的人负责关键路径：主循环、状态机、通信协议这类“别人都要依赖”的模块，不能交给进度最慢的人；
\item 每个组员都要有可演示的成果：避免一个人埋头写框架、其他人无事可做，最后验收时才发现框架方向错了；
\item 每天合并一次代码：哪怕功能不完整，也要让代码能编译、能跑，防止最后一天才发生“地狱合并”。
\end{enumerate}

在 AI 时代，还可以给每个组员配一个“AI 帮手”，但必须统一规则：用什么模型、遵守什么命名规范、改动哪些文件要提前声明。否则五个人各自让 AI 改五份代码，最后拼起来比手写还慢。

% ---------------------------------------------------------------------
\subsection{如何与实验室老师谈判}

很多同学不敢找老师谈延期，怕被批评；或者只会说“老师，我做不完了”。这两种方式都拿不到资源。老师真正想知道的不是“你有多惨”，而是“你还能交付什么”。

与老师谈判时建议按下面三步走：

\begin{enumerate}
\item 带着进度去，而不是带着困难去：把已经完成的功能现场演示一遍，哪怕只是一个能亮灯、能收发串口的半成品；
\item 给出二选一的具体方案：例如“方案 A：保留基础功能，砍掉数据存储，三天后可交付；方案 B：全部功能保留，但需要再给我一周调硬件”，让老师做选择题，而不是做判断题；
\item 把“完不成”的原因说成风险，而不是借口：硬件采购延迟、某个芯片手册与实物不符，这些是可验证的事实；说“太难了”则容易被理解为态度问题。
\end{enumerate}

记住：实验室老师的核心诉求通常是“这个项目能作为教学案例顺利收尾”，而不是为难你。只要你能证明自己在推进、并且给出一条可行的收尾路径，大部分老师愿意在范围或时间上让步。

% ---------------------------------------------------------------------
\subsection{为什么校园里一定要学会造“塑料轮子”}

这里举作者自己的例子。做毕业设计时，作者一开始认准了 ROS2，觉得它才是“正规轮子”，于是硬啃了三个月：安装环境、学习节点通信、调 EKF 参数……越到后期越焦虑，因为交期一天天逼近，而 ROS2 的工程量根本不是毕业设计剩下的时间能消化的。直到 DDL 真的压到面前，作者才决定先不学 ROS2，改造一个“塑料轮子”：硬件上用双核异构架构，STM32H743 负责运动控制与数据融合，K230 负责视觉识别；软件上放弃 Linux 和 ROS2，用 AI 生成的表达式解析器处理可配置逻辑，再把 EKF 换成基于动态打滑度检测的自适应里程计融合算法，两块芯片之间只靠一条高速串口通信。这套方案谈不上工程美感，但它只花了一周左右就跑通了，最后顺利交付。

为什么校园里一定要学会造塑料轮子？原因主要有几点：

\begin{enumerate}
\item 校园项目的 DDL 通常不可商量：课程设计、毕业设计、竞赛都有硬性截止日期，错过就只能等下一轮。先保证“能交付”，才有资格谈“做得好”；
\item 校园阶段没有“试错预算”：没有成熟的团队、售后和技术社区替你兜底，时间耗在框架学习上，往往意味着论文和答辩一起被拖垮；
\item 校园验收看的是“完整可演示”：老师评审的是功能能不能跑、思路能不能讲清楚，不是架构能不能支撑十年迭代。一个能跑通的塑料轮子，比一个永远学不完的正规框架更有价值；
\item 造塑料轮子本身就是学习：为了绕开 ROS2，你需要自己设计通信协议、写解析器、做传感器融合。这个过程逼你把“黑盒”拆开，反而比直接调用框架更接近底层原理。
\end{enumerate}

需要强调的是：造塑料轮子不等于拒绝学习 ROS2。作者后来再接触 ROS2 时发现，AI Agent 可以帮忙分析日志、批量调节参数，效率比当年手工死磕高得多。也就是说，塑料轮子是校园阶段保护自己的应急技能，不是拒绝成长的借口；等时间、工具都齐备了，再回头把塑料轮子换成正式轮子也不迟。

% ---------------------------------------------------------------------
\section{职场的紧急交期处理}

职场的交期和校园最大的不同，是它背后连着钱、合同和信任。一次延期可能影响整个项目排期，所以处理紧急交期不能只靠加班，更要靠“分析、取舍、谈判”。

% ---------------------------------------------------------------------
\subsection{分析哪些任务必须完成}

时间不够时，第一件事不是埋头写代码，而是先回答一个问题：如果明天就要交付，哪些事情不做就一定会失败？

判断标准可以从“验收视角”倒推：

\begin{itemize}
\item 老板或客户验收时会操作哪些功能？这些功能里的主流程就是必须完成的；
\item 哪些 Bug 会导致死机、数据丢失、无法通信？这些是致命缺陷，必须修；
\item 哪些只是“颜色不对、按钮位置不好看、代码可以更优雅”？这些属于优化，可以放在交期之后；
\item 哪些任务是别人提的“顺手做一下”，不在验收清单里？这类任务要敢于砍掉或明确推迟。
\end{itemize}

分析完之后，把任务分成“必须做”“应该做”“可以不做”三档，然后只承诺第一档的完成时间。需要特别提醒的是：砍任务之前一定要先和相关方说清楚，不能自己默默不做，否则交付时会被认为“漏功能”，而不是“主动取舍”。

% ---------------------------------------------------------------------
\subsection{使用别人的轮子还是自研轮子}

交期压力下，很多人第一反应是“找个现成库赶紧用”。但作者的经验是：轮子不是拿来就能跑的，选错轮子比没有轮子更危险。

如果这个轮子是业余爱好者维护的，或是靠程序员“用爱发电”开发的，即使是 MIT 协议，也要尽可能避免使用。原因很简单：用爱发电的项目，除了一些自律性很高的团体能够长期维护，大多数会随着维护者毕业、换工作、失去兴趣而停止更新。而你的产品一旦依赖它，就等于把交期押在了一个随时可能消失的人身上。

判断一个轮子能不能用，可以看两个日期：该项目的更新日期和目标框架的更新日期。如果两者相差很远，说明这个项目很可能已经跟不上框架的 API 变化，强行接入会出现大量“接口对接不上”的问题。此时迁移 API 的时间成本，往往不如自己写一个。

什么情况下应该自研？

\begin{itemize}
\item 需求非常简单，只有几十行，引入库反而要学它的封装；
\item 现有轮子的 API 与目标框架严重脱节，迁移成本高于重写成本；
\item 这个模块是产品的核心，未来要长期维护、深度定制，不能依赖一个黑盒；
\item 时间紧到没空读源码，但又必须保证每一行都可控。
\end{itemize}

这和第~\ref{ch:about-author}~章里“自研 Qt 绘图层替代 HelixScreen”是同一个逻辑：与其在看不透的封装层里反复试错，不如花一两天梳理需求，用确定的代码换取后面无数个不眠夜的安宁。

% ---------------------------------------------------------------------
\subsection{优化任务强度}

交期越紧，人越容易陷入“用战术上的勤奋掩盖战略上的懒惰”：笔记抄了一大堆、代码写了一大堆、Bug 也修了一大堆，但真正推动项目前进的只有那么几件事。

作者观察到一个现象：学霸学习时往往没有记录很多笔记，他们学习好的原因是会减少自己的记忆量，增加自己的思考时间。编程也一样——高手不是记得更多 API，而是知道哪些东西不必记。

优化任务强度可以从三个方向入手：

\begin{enumerate}
\item 减少记忆量：把命名规范写进 \texttt{AGENTS.md}、把常用操作写进 \texttt{SKILL.md}、把重复命令做成脚本，让 AI 和工具替你记；
\item 增加思考量：把省下来的时间用来想“这个功能真的需要吗”“这个方案三个月后还能改吗”，而不是用来背 API；
\item 降低切换成本：一次只专注一个模块，避免在菜单、通信、算法之间来回跳转，因为人的大脑切换上下文比编译器慢得多。
\end{enumerate}

用作者的话说：物理上的勤快要用来换取认知上的偷懒，而不是反过来。

% ---------------------------------------------------------------------
\subsection{同事间的甩锅技巧}

“甩锅技巧”听起来不太正面，但作者想表达的是另一层意思：紧急交期下，最怕的不是活多，而是责任边界不清——出了问题互相推诿，最后谁都在改代码，谁都没改对。

作者在第一家公司上班时，团队使用流水线式的 WBS 规划，每个人只负责自己那一小块。作者很快发现，产品里 UI 的模态框有很多种，每一种都长得差不多，却又各自独立实现，改一个需求就要把所有模态框都改一遍。于是他找到设计师同事，建议把所有模态框统一成一种组件。

这个建议的妙处在于：它不是在“甩锅”，而是在消灭“锅”本身。

\begin{itemize}
\item 模态框统一后，需求变更只需要改一个公共组件；
\item 每个同事的改动都基于同一套接口，联调时不会再出现“我这边没问题，是他那边格式不对”的争论；
\item 出问题时可以通过日志快速定位到具体模块，而不是在十几个相似实现里人肉搜索。
\end{itemize}

所以真正的“防甩锅技巧”是：

\begin{enumerate}
\item 把接口定义清楚：谁负责输入、谁负责输出、异常谁来处理，白纸黑字写下来；
\item 把重复的东西收敛：多种模态框合并成一种、多种消息格式统一成一种，公共部分越少，背锅点越少；
\item 让 Bug 能定位：每个模块有独立日志前缀和错误码，问题出现时先看日志再开会，而不是先开会对骂；
\item 重要决策留痕：邮件、会议纪要、代码注释里写清楚“为什么这么设计”，以后有争议时拿事实说话。
\end{enumerate}

% ---------------------------------------------------------------------
\subsection{如何与老板谈判}

程序员普遍不擅长谈判，不是嘴笨，而是总想用技术细节说服对方。老板不关心你用了什么设计模式，他只关心：什么时候能交付、要花多少钱、风险是什么。

作者遇到过这样一件事：设计师同事认为，要统一 UI 模态框就得找老板谈判，因为改动范围跨越多个模块。作者的建议是：直接去找老板，把核心需求说出来就好。

具体怎么说？

\begin{enumerate}
\item 先讲现状：现在系统里有十几种模态框，视觉不统一，每次加新功能都要重复开发；
\item 再讲方案：统一成一个公共组件，以后新增弹窗只需填参数；
\item 最后讲代价：需要两天重构时间，但之后每次改 UI 都能节省大量工时，长期收益远大于短期成本；
\item 给老板一个明确的问题：“是否同意用两天时间做这次统一？”
\end{enumerate}

因为程序员没几个说话利索的，这个时候最重要的不是展示口才，而是表述明确：把“我要什么、为什么、花多久、风险是什么”一次说清。老板最怕的不是你提要求，而是你提了一个他听不懂的要求。

谈判时还要守住一条底线：承诺过的交期一定要兑现，或者提前暴露风险。宁可一开始说“这个功能需要三天”，也不要最后说“其实昨天就该告诉你做不完”。信任一旦透支，下一次再紧急的交期，就没有人愿意听你谈判了。

% ---------------------------------------------------------------------

总结本章：面对紧急交期，AI 负责提高“体力上限”，自研还是用轮子决定“风险下限”，而谈判能力决定“交付底线”。三者合在一起，就是作者在第~\ref{ch:about-author}~章里反复强调的那句话——活着通关才是硬道理。
